
\documentclass{article}
\usepackage{colortbl}
\usepackage{makecell}
\usepackage{multirow}
\usepackage{supertabular}

\begin{document}

\newcounter{utterance}

\centering \large Interaction Transcript for game `adventuregame', experiment `home\_deliver\_three\_basic\_hard\_invlimittwo', episode 1 with reka{-}flash{-}3.1{-}t0.0.
\vspace{24pt}

{ \footnotesize  \setcounter{utterance}{1}
\setlength{\tabcolsep}{0pt}
\begin{supertabular}{c@{$\;$}|p{.15\linewidth}@{}p{.15\linewidth}p{.15\linewidth}p{.15\linewidth}p{.15\linewidth}p{.15\linewidth}}
    \# & \multicolumn{2}{c}{Player} && \multicolumn{2}{c}{Game Master} \\
    \hline

    \theutterance \stepcounter{utterance}  
    & & & \multicolumn{4}{p{0.6\linewidth}}{
        \cellcolor[rgb]{0.9,0.9,0.9}{
            \makecell[{{p{\linewidth}}}]{
                \texttt{\tiny{[P1$\langle$GM]}}
                \texttt{You are playing a text adventure game. I will describe what you can perceive in the game. You write the single action you want to take in the game starting with >. Only reply with an action.} \\
\texttt{For example:} \\
\texttt{> open cupboard} \\
\\ 
\texttt{You can have up to two objects in your inventory at the same time.} \\
\\ 
\texttt{Your goal for this game is: Put the sandwich in the freezer, the plate on the shelf and the mop on the table.} \\
\\ 
\texttt{Once you have achieved your goal, write "> done" to end the game.} \\
\\ 
\texttt{You are in a bedroom now. There are a pillow, a wardrobe and a bed. The pillow is in the wardrobe. The wardrobe is closed. There is a passage to a hallway here.} \\
            }
        }
    }
     \\ \\

    \theutterance \stepcounter{utterance}  
    & \multicolumn{4}{p{0.6\linewidth}}{
        \cellcolor[rgb]{0.561,0.737,0.561}{
            \makecell[{{p{\linewidth}}}]{
                \texttt{\tiny{[P1$\rangle$GM]}}
                \texttt{<reasoning>} \\
\texttt{Let's see. The goal is to put the sandwich in the freezer, the plate on the shelf, and the mop on the table. But right now, I'm in the bedroom. The description says there's a pillow, a wardrobe, and a bed. The pillow is in the wardrobe, and the wardrobe is closed. There's a passage to the hallway.} \\
\\ 
\texttt{Wait, but the items needed for the goals are sandwich, plate, mop. But the bedroom has pillow, wardrobe, bed. Maybe I need to go to the hallway first? Because the passage is there. Let me check the initial description again. The user says "You are in a bedroom now. There are a pillow, a wardrobe and a bed. The pillow is in the wardrobe. The wardrobe is closed. There is a passage to a hallway here." So the hallway is accessible via the passage. But the items needed (sandwich, plate, mop) aren't mentioned here. So maybe those items are in the hallway or other rooms connected to it?} \\
\\ 
\texttt{But the user hasn't described the hallway yet. So the first step is probably to go to the hallway. Let me think. The initial state is the bedroom. The user needs to move to the hallway to find the other items. So the first action should be to go to the hallway. Let me check the example: the example was "open cupboard" as an action. So moving is allowed. So the action would be "> go to hallway" or "> go hallway". But the user might expect "go to hallway" or "enter hallway". Wait, the example uses "open cupboard" which is a verb{-}object. So maybe "go to hallway" is the correct format. Let me check the initial problem statement. The user says "You write the single action you want to take in the game starting with >. Only reply with an action." So the action is a command like "open cupboard" or "go to hallway". So the first step is to go to the hallway. So the action would be "> go to hallway".} \\
\\ 
\texttt{Once in the hallway, what's there? The user hasn't described it yet, but maybe the hallway has other rooms or items. But since the goal items are sandwich, plate, mop, perhaps those are in the hallway or other connected areas. Alternatively, maybe the bedroom has those items but they're not mentioned. Wait, the initial description says "You are in a bedroom now.} \\
            }
        }
    }
    & & \\ \\

    \theutterance \stepcounter{utterance}  
    & & & \multicolumn{2}{p{0.3\linewidth}}{
        \cellcolor[rgb]{0.9,0.9,0.9}{
            \makecell[{{p{\linewidth}}}]{
                \texttt{\tiny{[GM$|$GM]}}
                \texttt{Round: 0} \\
            }
        }
    }
    & & \\ \\

    \theutterance \stepcounter{utterance}  
    & & & \multicolumn{2}{p{0.3\linewidth}}{
        \cellcolor[rgb]{0.9,0.9,0.9}{
            \makecell[{{p{\linewidth}}}]{
                \texttt{\tiny{[GM$|$GM]}}
                \texttt{command\_tag\_missing} \\
            }
        }
    }
    & & \\ \\

    \theutterance \stepcounter{utterance}  
    & & & \multicolumn{2}{p{0.3\linewidth}}{
        \cellcolor[rgb]{0.9,0.9,0.9}{
            \makecell[{{p{\linewidth}}}]{
                \texttt{\tiny{[GM$|$GM]}}
                {'goal_states_achieved': [], 'game_successfully_finished': False}
            }
        }
    }
    & & \\ \\

    \theutterance \stepcounter{utterance}  
    & & & \multicolumn{2}{p{0.3\linewidth}}{
        \cellcolor[rgb]{0.9,0.9,0.9}{
            \makecell[{{p{\linewidth}}}]{
                \texttt{\tiny{[GM$|$GM]}}
                \texttt{command\_tag\_missing} \\
            }
        }
    }
    & & \\ \\

\end{supertabular}
}

\end{document}
